% =====================================================================
% Part I — TotD cluster A (grupos de permutaciones). Edición española.
% Autores: Carles Marín (Claude, Anthropic, como asistente IA).
% Traducción paralela de neumann-separation-lean.tex.
% =====================================================================
\documentclass[11pt]{article}
\usepackage[T1]{fontenc}
\usepackage[spanish,es-noshorthands]{babel}
\usepackage{lmodern}
\usepackage{amsmath,amssymb,amsthm,booktabs,graphicx,hyperref}
\usepackage{xurl}
\usepackage[table]{xcolor}
\usepackage{pifont}
\usepackage{multirow}
\usepackage{tikz}
\usetikzlibrary{arrows.meta,positioning,calc,shapes.geometric}
\definecolor{sepblue}{HTML}{1B4F8C}
\definecolor{sepband}{HTML}{DCE3EB}
\definecolor{sepamber}{HTML}{E08A1E}
\definecolor{sepred}{HTML}{C0392B}
\definecolor{voiceA}{HTML}{1B6F8C}
\definecolor{voiceB}{HTML}{E08A1E}
\definecolor{voiceC}{HTML}{8C2D6F}
\definecolor{restgrey}{HTML}{B9C2CC}
\definecolor{dTwo}{HTML}{DCE7EB}
\definecolor{dThree}{HTML}{F4C46B}
\definecolor{dFour}{HTML}{D9663A}
\usepackage[margin=1in]{geometry}
\usepackage{microtype}
\setlength{\emergencystretch}{2em}
\newtheorem{theorem}{Teorema}
\newtheorem{lemma}{Lema}
\newtheorem{definition}{Definición}
\newcommand{\lean}[1]{\texttt{#1}}
\newcommand{\Z}{\mathbb{Z}}
\newcommand{\R}{\mathbb{R}}
\newcommand{\N}{\mathbb{N}}
\newcommand{\GL}{\mathrm{GL}}

\title{\bf Cuando la separación falla, aparece un canon\\[3pt]
  \large\normalfont Una identidad de conteo tras la separación, la sincronización
  y la ruptura de simetría en grupos de permutaciones, formalizada en Lean~4
  (Parte~I)}
\author{Carles Mar\'in\thanks{Formalizado con asistencia de IA (Claude,
  Anthropic); las matemáticas y todas las afirmaciones son responsabilidad del
  autor.}\\[3pt]
  \normalsize Investigador independiente\\
  \normalsize \texttt{karlesmarin@gmail.com}}
\date{Julio 2026}

\begin{document}
\maketitle

\begin{abstract}
Una identidad aritmética recorre este artículo: el promedio, sobre un grupo finito
$G$ que actúa en $\Omega$, de cuántos puntos comparte una traslación $gA$ con un
conjunto $B$ es exactamente $|A|\,|B|/|\Omega|$. Leída de un modo demuestra el lema
de separación de P.~M. Neumann; leída de otro fuerza los teoremas de uniformidad de
la jerarquía de sincronización de P.~J. Cameron; y en la frontera
$|A|\,|B|=|\Omega|$ produce una dicotomía---o algún $g$ separa $A$ de $B$, o
\emph{toda} traslación corta a $B$ en exactamente un punto. Ese caso extremal es una
factorización exacta del grupo, y sobre $\Z/n$ un canon rítmico por teselación: la
frontera de la teoría de separación y la teoría de los cánones musicales son el
mismo objeto, en dos literaturas que no se citan entre sí. Formalizo este
conglomerado en Lean~4 / Mathlib, junto con el \emph{número distinguidor} de una
acción de grupo y el valor del plano de Fano $D(\GL_3(\mathbb{F}_2))=4$. La
matemática es clásica; lo nuevo aquí es la forma órbita-local de la dicotomía, la
identificación de su caso extremal con los cánones rítmicos y---hasta donde
sé---la primera formalización de cualquiera de estos en un asistente de pruebas.
Todo es \lean{sorry}-free, dependiendo solo de los tres axiomas estándar: una perla,
y la primera parte de una serie planeada.
\end{abstract}

% ---------------------------------------------------------------
\section{Un argumento, varias caras}
\label{sec:intro}

Sea un grupo finito $G$ actuando en un conjunto $\Omega$, y tomemos dos subconjuntos
finitos $A,B\subseteq\Omega$. Aplicar $g\in G$ desliza $A$ a su traslación $gA$.
Casi todo en este artículo se sigue de contar, de dos maneras, las ternas
$(a,b,g)$ con $a\in A$, $b\in B$ y $g\cdot a=b$. Cuando $G$ actúa transitivamente en
$\Omega$, sumar sobre el grupo y despejar denominadores da la identidad
\[
  \boxed{\;\Bigl(\sum_{g\in G}|gA\cap B|\Bigr)\cdot|\Omega|
  \;=\;|A|\cdot|B|\cdot|G|\;}
\]
---el tamaño promedio de $gA\cap B$ es $|A|\,|B|/|\Omega|$; una forma local a las
órbitas vale para cualquier acción, y es la que necesita la dicotomía de la
Sección~\ref{sec:dichotomy}. Parece una observación de
contabilidad. Es, de hecho, el motor de un conglomerado sorprendentemente amplio de
teoremas, y el único mensaje de este artículo es que son un argumento visto desde
distintos lados (Figura~\ref{fig:map}).

\begin{figure}[t]
  \centering
  \begin{tikzpicture}[
      font=\small,
      box/.style={rounded corners=2pt,draw=sepblue,line width=0.9pt,fill=sepband,
                  align=center,inner sep=5pt,minimum height=8mm},
      hero/.style={rounded corners=2pt,draw=sepblue,line width=1.3pt,fill=sepblue,
                   text=white,align=center,inner sep=6pt},
      star/.style={rounded corners=2pt,draw=voiceC,line width=1.1pt,fill=voiceC!12,
                   align=center,inner sep=5pt},
      dual/.style={rounded corners=2pt,draw=restgrey,line width=0.9pt,fill=white,
                   align=center,inner sep=5pt,text=black!75},
      fl/.style={-{Stealth[length=2.4mm]},sepblue,line width=0.9pt}]
    \node[hero] (id) at (0,5.2) {\textbf{Una identidad de conteo}\\[1pt]
      $\bigl(\sum_g|gA\cap B|\bigr)|\Omega|=|A||B||G|$};
    \node[box] (sep) at (-2.6,3.5) {Separación\\{\footnotesize (Neumann)}};
    \node[box] (uni) at (2.6,3.5) {Uniformidad\\{\footnotesize (Cameron 8--9)}};
    \node[box] (dic) at (0,2.0) {Dicotomía de separación\\{\footnotesize en $|A||B|=|\Omega|$}};
    \node[box] (fac) at (0,0.7) {Factorización exacta\\{\footnotesize $G=B\cdot A^{-1}$}};
    \node[star] (can) at (0,-0.7) {\textbf{Canon rítmico}\\{\footnotesize $A\oplus B=\Z/n$}};
    \node[dual] (dist) at (4.55,1.35) {\emph{La lente dual}\\ Número distinguidor\\{\footnotesize Fano $D=4$}};
    \draw[fl] (id) -- (sep); \draw[fl] (id) -- (uni);
    \draw[fl] (sep) -- (dic); \draw[fl] (uni) -- (dic);
    \draw[fl] (dic) -- (fac); \draw[fl] (fac) -- (can);
    \draw[-{Stealth[length=2.4mm]},restgrey,line width=0.9pt,dashed]
      (id.east) to[bend left=22] node[right,midway,text=black!55,font=\scriptsize]{kit órbita--}
      node[right,midway,yshift=-3.2mm,text=black!55,font=\scriptsize]{estabilizador} (dist.north);
  \end{tikzpicture}
  \caption{El artículo en una imagen. Una única identidad de conteo fuerza tanto la
  separación de Neumann como los teoremas de uniformidad de Cameron; donde se
  encuentran, en el umbral $|A||B|=|\Omega|$, produce una dicotomía cuyo caso no
  separable es una factorización exacta del grupo---sobre $\Z/n$, un canon rítmico
  por teselación. El número distinguidor (la excepción de Fano $D=4$) es una
  pregunta paralela sobre romper la simetría, que comparte el mismo kit
  órbita--estabilizador.}
  \label{fig:map}
\end{figure}

\paragraph{Qué es nuevo, en tres registros.} Como este artículo se sitúa entre un
survey y una formalización, separo sus afirmaciones con nitidez.
\emph{Matemáticamente}, el contenido nuevo es una forma afinada, órbita-local, de la
dicotomía de separación (Sección~\ref{sec:dichotomy}), y la observación de que su
caso extremal, sobre un grupo cíclico, es exactamente el canon rítmico por
teselación de Vuza~\cite{Vuza1991} y Coven--Meyerowitz~\cite{CovenMeyerowitz1999}.
\emph{En un asistente de pruebas}: hasta donde sé, tras una búsqueda bibliográfica
extensa por los sistemas principales, esta es la primera formalización del lema de
separación, de la jerarquía de sincronización de Cameron y del número distinguidor
de una acción de grupo. \emph{Conceptualmente}, la contribución es hacer explícito
el motor compartido: separación, uniformidad, la dicotomía extremal y la
factorización se reducen todas, por nombre y en el grafo de dependencias, a la única
identidad de arriba. El Apéndice~\ref{app:results} lista los resultados formales con
sus nombres Lean; la matemática es clásica en todo momento, y nombro a sus autores
según avanzo.

% Tabla del cuerpo (Tabla de contrastes) + macros de estado (edición ES).
% La tabla de resultados formales vive en el apéndice del fichero principal.
\newcommand{\stproven}{\textcolor{sepblue}{\ding{51}}\,probado}
\newcommand{\stfuture}{\textcolor{restgrey}{$\circ$}\,futuro}
\newcommand{\stdef}{\textcolor{sepblue}{$\bullet$}\,definido}

% ---------- Tabla de contrastes computacionales ----------
\begin{table}[t]
  \centering
  \small
  \arrayrulecolor{sepblue}
  \caption{Contrastes computacionales (GAP / Sage, exactos). \emph{Izquierda:} la
  dicotomía de separación, verificada sin contraejemplo en todo grupo transitivo de
  grado pequeño, más el testigo de canon en $\Z/12$. \emph{Derecha:} una
  recomputación independiente de la tabla distinguidora de Chan
  $D(\GL_n(\mathbb{F}_q)\curvearrowright\mathbb{F}_q^n)$ para $n>2$, $q\le n{+}1$,
  coincidente con Klavžar--Wong--Zhu; el color de celda codifica $D$ (pálido $D{=}2$
  $\to$ intenso $D{=}4$), así que las dos excepciones destacan de un vistazo.}
  \label{tab:checks}
  \renewcommand{\arraystretch}{1.2}
  \begin{minipage}[t]{0.46\textwidth}
    \centering
    \begin{tabular}{@{}ll@{}}
      \toprule
      \rowcolor{sepblue}
      \textcolor{white}{\textbf{Contraste}} & \textcolor{white}{\textbf{Resultado}} \\
      \midrule
      \rowcolor{sepband}
      barrido dicotomía & \multirow{1}{*}{0 contraejemplos} \\
      \small 120 grps.\ transitivos, gr.\ 2--9 & \\
      canon $\Z/12$ $\{0,3,6,9\}{\oplus}\{0,1,2\}$ & tesela (extremal) \\
      \bottomrule
    \end{tabular}
  \end{minipage}\hfill
  \begin{minipage}[t]{0.5\textwidth}
    \centering
    \setlength{\tabcolsep}{5pt}
    \begin{tabular}{@{}c|cccc@{}}
      \multicolumn{1}{@{}c}{\scriptsize $n\backslash q$} & \scriptsize 2 & \scriptsize 3 & \scriptsize 4 & \scriptsize 5 \\
      \hline
      \scriptsize 3 & \cellcolor{dFour}\textbf{4} & \cellcolor{dTwo}2 & \cellcolor{dTwo}2 & --- \\
      \scriptsize 4 & \cellcolor{dThree}\textbf{3} & \cellcolor{dTwo}2 & \cellcolor{dTwo}2 & \cellcolor{dTwo}2 \\
      \scriptsize 5 & \cellcolor{dTwo}2 & \cellcolor{dTwo}2 & \cellcolor{dTwo}2 & \cellcolor{dTwo}2 \\
      \scriptsize 6 & \cellcolor{dTwo}2 & \cellcolor{dTwo}2 & \cellcolor{dTwo}2 & \cellcolor{dTwo}2 \\
    \end{tabular}
    \\[3pt]
    {\scriptsize\colorbox{dTwo}{\,$D{=}2$\,}\ \colorbox{dThree}{\,$D{=}3$\,}\ \colorbox{dFour}{\textcolor{white}{\,$D{=}4$\,}}\quad(``---'': $q$ no es potencia de primo)}
  \end{minipage}
\end{table}


% ---------------------------------------------------------------
\section{Separar un conjunto de su imagen}
\label{sec:setting}

La cara más antigua de la identidad es una pregunta de respuesta intuitiva:
\begin{quote}\itshape
  ¿cuándo podemos hallar siempre un $g$ que aparte $A$ enteramente de $B$---con
  $gA\cap B=\emptyset$?
\end{quote}
Si las órbitas son grandes hay sitio para deslizar $A$ lejos de $B$. Neumann lo
precisó. En el caso finito transitivo el umbral limpio es $|A|\cdot|B|<|\Omega|$: por
debajo, la separación siempre es posible (Figura~\ref{fig:sep}); la identidad de
conteo es toda la prueba, pues si ningún $g$ separara, cada término $|gA\cap B|\ge1$,
forzando la suma a al menos $|G|$ y el promedio a al menos $1$, es decir
$|A|\,|B|\ge|\Omega|$.

\begin{figure}[t]
  \centering
  \begin{tikzpicture}[scale=0.88,
      pt/.style={circle,draw=restgrey,line width=0.6pt,fill=white,minimum size=10pt,inner sep=0pt},
      inA/.style={circle,draw=sepblue,line width=1pt,fill=sepblue,minimum size=10pt,inner sep=0pt},
      inB/.style={circle,draw=sepred,line width=1.4pt,fill=white,minimum size=15pt,inner sep=0pt}]
    \begin{scope}[shift={(0,0)}]
      \foreach \i in {0,...,9} { \coordinate (l\i) at ({90-36*\i}:1.7); }
      \draw[restgrey!60] (0,0) circle (1.7);
      \foreach \i in {0,...,9} { \node[pt] at (l\i) {}; }
      \foreach \i in {5,6} { \node[inB] at (l\i) {}; }
      \foreach \i in {4,5,6} { \node[inA] at (l\i) {}; }
      \node at (0,-2.25) {\small \textbf{antes:} $A\cap B\neq\emptyset$};
    \end{scope}
    \draw[-{Stealth[length=3.5mm]},sepamber,line width=1.3pt]
      (2.15,0.15) to[bend left=18] node[above,midway,text=sepamber]{\small rota por $g$} (3.55,0.15);
    \begin{scope}[shift={(5.7,0)}]
      \foreach \i in {0,...,9} { \coordinate (r\i) at ({90-36*\i}:1.7); }
      \draw[restgrey!60] (0,0) circle (1.7);
      \foreach \i in {0,...,9} { \node[pt] at (r\i) {}; }
      \foreach \i in {5,6} { \node[inB] at (r\i) {}; }
      \foreach \i in {0,1,2} { \node[inA] at (r\i) {}; }
      \node at (0,-2.25) {\small \textbf{después:} $gA\cap B=\emptyset$};
    \end{scope}
    \begin{scope}[shift={(-1.3,-3.15)}]
      \node[inA,scale=0.8] at (0,0) {}; \node[right=1pt,anchor=west] at (0.15,0) {\scriptsize $A$ ($|A|{=}3$)};
      \node[inB,scale=0.7] at (2.4,0) {}; \node[right=1pt,anchor=west] at (2.6,0) {\scriptsize $B$ ($|B|{=}2$)};
      \node[anchor=west] at (4.6,0) {\scriptsize $|A||B|=6<10=|\Omega|$};
    \end{scope}
  \end{tikzpicture}
  \caption{Separación en un decágono (acción regular de $\Z/10$). Como
  $|A|\cdot|B|<|\Omega|$, alguna rotación desliza $A$ (azul) fuera de $B$ (anillos
  rojos). El umbral $|A||B|=|\Omega|$, donde la separación puede fallar, es donde la
  historia gira (Sección~\ref{sec:dichotomy}).}
  \label{fig:sep}
\end{figure}

Una sutileza merece una frase, porque una paráfrasis popular la yerra. El umbral es
\emph{estricto}. La fuente cuantitativa---Birch, Burns, Oates Macdonald y
Neumann~\cite{BBMN1976}---separa bajo la hipótesis de que toda órbita tenga
\emph{más de} $|A|\,|B|$ elementos, y observa que $|A|\,|B|+1$ es óptimo: un grupo
transitivo de grado $|A|\,|B|$ construido con $|A|$ bloques de tamaño $|B|$ tiene un
par no separable. Debilitar «más de» a «al menos» hace falsa la afirmación (el grupo
trivial en un punto nunca separa), y el caso frontera $|A|\,|B|=|\Omega|$---donde el
lema estricto calla---es exactamente donde vive la Sección~\ref{sec:dichotomy}.

Para $\Omega$ infinito la hipótesis correcta es sobre órbitas, y la forma afinada
pide solo que la órbita de cada punto de $A$ sea infinita. La prueba es la dual
placentera de la finita: si ningún $g$ separara, los cosets $\{g\mid g\cdot a=b\}$
cubrirían $G$, así que por el lema de recubrimiento de B.~H.
Neumann~\cite{BHNeumann1954} algún estabilizador de punto tendría índice finito---pero
su índice es el tamaño (infinito) de la órbita. Mathlib ya lleva el lema de
recubrimiento de B.~H. Neumann; el envoltorio para acciones de grupo que le faltaba
se aporta aquí.

% ---------------------------------------------------------------
\section{Uniformidad: el mismo promedio, leído como igualdad}
\label{sec:engine}

Las lecciones de sincronización de Cameron~\cite{CameronSync} convierten la identidad
en una familia de enunciados de rigidez. Si toda intersección-traslación tiene por
casualidad el \emph{mismo} tamaño $\ell$, la identidad da de inmediato
$|A|\,|B|=\ell\,|\Omega|$; su instancia $\ell=1$ dice que una partición
sección-regular es uniforme, todas las partes del mismo tamaño. Más fino, y el hecho
que impulsa la sección siguiente: si todas las intersecciones son meramente $\ge\ell$
(o todas $\le\ell$) mientras $|A|\,|B|=\ell\,|\Omega|$, entonces cada una vale
exactamente $\ell$.

El punto notable merece decirse en palabras llanas. Un promedio es un enunciado sobre
una \emph{suma}; aquí baja y fija cada término individual. Como la media de los
$|gA\cap B|$ está forzada a ser exactamente $\ell$, ninguna traslación puede quedar
por encima ni por debajo sin que otra compense---así que en cuanto una cota unilateral
retira esa libertad, el promedio dicta cada término por sí solo. Es un promedio que
legisla puntualmente.

Soy cuidadoso con lo que es esta «unificación». No probé un meta-teorema del que la
separación, la uniformidad y la dicotomía sean corolarios; más bien, las tres llaman
al mismo conteo de ternas. La unidad es real pero modesta, y vive en el grafo de
dependencias, no en un enunciado grandioso---que es precisamente el tipo de afirmación
que una formalización deja comprobar a un lector en vez de aceptar por fe.

% ---------------------------------------------------------------
\section{Cuando la separación falla: una dicotomía, y un canon}
\label{sec:dichotomy}

Volvamos a la frontera $|A|\cdot|B|=|\Omega|$, donde el lema estricto de separación no
dice nada. Aquí el motor fuerza una bifurcación limpia.

\begin{theorem}[Dicotomía de separación, forma órbita-local]
Sea $G$ finito actuando en $\Omega$, y $A,B$ finitos no vacíos con toda órbita que
corta a la vez $A$ y $B$ de tamaño $\ge|A|\cdot|B|$. Entonces o bien
\begin{enumerate}
\item[\rm(1)] algún $g$ separa $A$ de $B$, o bien
\item[\rm(2)] \emph{(extremal)} $|gA\cap B|=1$ para \emph{todo} $g\in G$, y $A$ y $B$
  yacen en una única órbita común de tamaño exactamente $|A|\cdot|B|$.
\end{enumerate}
\end{theorem}

La prueba es el promedio una vez más: si ningún $g$ separa, cada término es $\ge1$; la
identidad fija su media en $1$; así que por el estrujón puntual de la sección anterior
cada $|gA\cap B|$ vale exactamente $1$, y el caso de igualdad colapsa $A\cup B$ en una
órbita del tamaño extremal. Estrictamente por encima del umbral se obtiene separación
sin más.

Enuncio la honestidad con claridad, porque el núcleo transitivo de la bifurcación no
es mío. Que $|A|\,|B|=|\Omega|$ fuerce o bien separación o bien
todas-las-intersecciones-uno es, en esencia literalmente, la definición de un
\emph{grupo separador} en Araújo--Cameron--Steinberg~\cite{ACS2017}. Lo que reclamo es
la forma órbita-local, débil, no transitiva de arriba---la hipótesis restringe solo las
órbitas que cortan ambos conjuntos, y la conclusión \emph{produce} la única órbita
común---y su formalización.

\paragraph{El caso extremal es una factorización.} ¿Cuáles son los pares extremales?
Para la acción regular de $G$ sobre sí mismo, «$|gA\cap B|=1$ para todo $g$» dice
exactamente que todo elemento del grupo es de forma única $b\,a^{-1}$ con $a\in A$,
$b\in B$: una \emph{factorización exacta} $G=B\cdot A^{-1}$. Escrita aditivamente para
$G=\Z/n$ hace del mapa diferencia $(a,b)\mapsto b-a$ una biyección; equivalentemente,
como $|A|\,|B|=n$, las traslaciones $\{A+b : b\in B\}$ particionan $\Z/n$---la
teselación $A\oplus B=\Z/n$. Tal teselación es un \emph{canon rítmico}: un ritmo $A$ y
un conjunto de entradas $B$ cuyas copias desplazadas cubren cada pulso una vez, sin
hueco ni choque (Figura~\ref{fig:canon}). Un par extremal no separable \emph{es} un
canon.

\begin{figure}[t]
  \centering
  \begin{tikzpicture}[scale=1.0]
    \begin{scope}[shift={(-3.4,0)}]
      \draw[restgrey!55,line width=0.8pt] (0,0) circle (2.15);
      \foreach \i in {0,...,11} { \coordinate (c\i) at ({90-30*\i}:2.15); }
      \foreach \i in {0,3,6,9}  { \fill[voiceA] (c\i) circle (7pt); }
      \foreach \i in {1,4,7,10} { \fill[voiceB] (c\i) circle (7pt); }
      \foreach \i in {2,5,8,11} { \fill[voiceC] (c\i) circle (7pt); }
      \foreach \i in {0,...,11} { \node[white,font=\scriptsize\bfseries] at (c\i) {\i}; }
      \node at (0,-2.7) {\small cada pulso cubierto \emph{exactamente una vez}};
    \end{scope}
    \begin{scope}[shift={(0.2,-1.55)}]
      \foreach \t in {0,...,12} { \draw[restgrey!35,line width=0.4pt] (\t*0.42,0) -- (\t*0.42,2.1); }
      \foreach \r in {0,1,2,3} { \draw[restgrey!35,line width=0.4pt] (0,\r*0.7) -- (5.04,\r*0.7); }
      \foreach \t in {0,3,6,9} { \node[font=\tiny,text=restgrey!80] at (\t*0.42+0.21,-0.28) {\t}; }
      \foreach \t in {0,3,6,9}  { \fill[voiceA] (\t*0.42+0.05,1.45) rectangle (\t*0.42+0.37,2.05); }
      \foreach \t in {1,4,7,10} { \fill[voiceB] (\t*0.42+0.05,0.75) rectangle (\t*0.42+0.37,1.35); }
      \foreach \t in {2,5,8,11} { \fill[voiceC] (\t*0.42+0.05,0.05) rectangle (\t*0.42+0.37,0.65); }
      \node[left,font=\scriptsize,text=voiceA] at (-0.1,1.75) {v.\,1};
      \node[left,font=\scriptsize,text=voiceB] at (-0.1,1.05) {v.\,2};
      \node[left,font=\scriptsize,text=voiceC] at (-0.1,0.35) {v.\,3};
      \node[font=\scriptsize] at (2.5,2.5) {\small entradas $B=\{0,1,2\}$, cada voz toca $A=\{0,3,6,9\}$};
    \end{scope}
  \end{tikzpicture}
  \caption{El par extremal (no separable) $A=\{0,3,6,9\}$, $B=\{0,1,2\}$ en $\Z/12$ es
  un canon rítmico. Izquierda: las tres traslaciones $A$ (turquesa), $A{+}1$ (ámbar),
  $A{+}2$ (violeta) particionan los doce pulsos---sin hueco, sin solapamiento.
  Derecha: como música, tres voces entran a un pulso de distancia, cada una tocando el
  ritmo $A$. Una versión sonificada acompaña a la nota musical compañera.}
  \label{fig:canon}
\end{figure}

Señalo los vecinos cercanos para que un revisor no tenga que hacerlo. La idea «no
separación $=$ factorización» aparece, para acciones diagonales, en
Araújo--Cameron--Steinberg~\cite{ACS2017} (factorizaciones perfectas de grupos
simples), y para el caso $A=B$ en Barbieri et~al.~\cite{Barbieri2025} (conjuntos
diferencia y planos proyectivos). Lo que parece no señalado es la identificación del
caso extremal regular de dos conjuntos sobre $\Z/n$ con los cánones rítmicos por
teselación: la literatura de separación y la de teselación (Vuza~\cite{Vuza1991},
Coven--Meyerowitz~\cite{CovenMeyerowitz1999}, Amiot~\cite{Amiot2016}) no se citan
mutuamente. La conexión es nueva como \emph{observación}, no como matemática; su valor
es que ahora es un lema en vez de una analogía---que es para lo que sirve un asistente
de pruebas.

% ---------------------------------------------------------------
\section{La pregunta dual: ¿cuánta simetría hay que romper?}
\label{sec:distinguishing}

Hasta aquí el motor de conteo ha medido cómo un grupo \emph{mueve} conjuntos---cuánto
puede empujar $A$ fuera de $B$. Una cantidad espejo mide cuán difícil es \emph{fijar}
la acción: cuánto etiquetado hace falta para eliminar su simetría por completo. Veo a
las dos como extremos opuestos de un mismo programa---la separación añade sitio hasta
que el grupo puede mover un conjunto a un lado; la distinción añade colores hasta que
el grupo no puede mover nada---y una única desigualdad probada abajo las une: una base
de tamaño $b$ (un conjunto cuyo estabilizador puntual es trivial) da $D\le b+1$, la
cota que ata la distinción a la teoría de bases y matroides (IBIS) que forma la
espina dorsal de la Parte~II. Por eso el tema pertenece aquí, como la cara dual del
programa de conteo y no como un apéndice.

Coloreamos los puntos de $\Omega$; el coloreo es \emph{distinguidor} si la identidad es
el único elemento de $G$ que lo preserva, y el \emph{número distinguidor}
$D(G\curvearrowright\Omega)$ es el mínimo de colores que basta
(Albertson--Collins~\cite{AlbertsonCollins1996}; Tymoczko~\cite{Tymoczko2004} para
acciones generales). Mathlib no tenía noción de ello; la definición, su monotonía y la
cota de base $D\le b+1$ se aportan aquí.

\begin{quote}\itshape
  Para $\GL_n(\mathbb{F}_q)$ actuando en los vectores no nulos de $\mathbb{F}_q^n$,
  ¿cuántos colores hacen falta?
\end{quote}
Melody Chan preguntó exactamente esto~\cite{Chan2006}; la respuesta es
Klavžar--Wong--Zhu~\cite{KWZ2006}: dos colores bastan siempre, con cuatro excepciones,
de las que el máximo único es $D(\GL_3(\mathbb{F}_2))=4$. Ese grupo, de orden $168$, es
el grupo de colineaciones del plano de Fano en sus siete puntos, y el valor excepcional
es el que verifico (Figura~\ref{fig:fano}).

La sorpresa es el tamaño del salto. Para toda acción lineal mayor dos colores rompen
toda la simetría; aquí, en solo siete puntos, ni siquiera \emph{tres} colores
pueden---cada uno de los $3^7=2187$ tres-coloreos es preservado por alguna
colineación---mientras que un cuatro-coloreo concreto sí. La cota inferior es la mitad
exhaustiva: veintiuna colineaciones bastan para fijar entre ellas cada tres-coloreo,
comprobado por el núcleo. La cota superior es pura geometría: tres puntos llevan
colores únicos y quedan fijados, tres rectas por ellos fijan tres más, y la
inyectividad fija el último. El plano de Fano es lo bastante pequeño para verlo y lo
justo de grande para ser rígido.

\begin{figure}[t]
  \centering
  \begin{tikzpicture}[scale=1.5,
      fp/.style={circle,draw=black!55,line width=0.7pt,minimum size=13pt,inner sep=0pt,font=\scriptsize\bfseries}]
    \coordinate (P0) at (90:1.5);
    \coordinate (P1) at (210:1.5);
    \coordinate (P2) at (330:1.5);
    \coordinate (P3) at ($(P0)!0.5!(P1)$);
    \coordinate (P4) at ($(P1)!0.5!(P2)$);
    \coordinate (P5) at ($(P2)!0.5!(P0)$);
    \coordinate (P6) at (0,0);
    \draw[sepblue,line width=0.9pt] (P0)--(P1)--(P2)--cycle;
    \draw[sepblue,line width=0.9pt] (P0)--(P4); \draw[sepblue,line width=0.9pt] (P1)--(P5); \draw[sepblue,line width=0.9pt] (P2)--(P3);
    \draw[sepblue,line width=0.9pt] (P6) circle (0.75);
    \node[fp,fill=restgrey,text=white] at (P0) {0};
    \node[fp,fill=restgrey,text=white] at (P1) {1};
    \node[fp,fill=restgrey,text=white] at (P2) {2};
    \node[fp,fill=restgrey,text=white] at (P3) {3};
    \node[fp,fill=voiceB,text=white] at (P4) {4};
    \node[fp,fill=voiceC,text=white] at (P5) {5};
    \node[fp,fill=sepred,text=white]  at (P6) {6};
    \begin{scope}[shift={(2.35,1.15)},every node/.style={anchor=west,font=\scriptsize}]
      \fill[restgrey] (0,0) circle (3pt);   \node at (0.18,0) {color 0: $\{0,1,2,3\}$};
      \fill[voiceB]  (0,-0.42) circle (3pt);\node at (0.18,-0.42) {color 1: $\{4\}$};
      \fill[voiceC]  (0,-0.84) circle (3pt);\node at (0.18,-0.84) {color 2: $\{5\}$};
      \fill[sepred]  (0,-1.26) circle (3pt);\node at (0.18,-1.26) {color 3: $\{6\}$};
    \end{scope}
    \node[font=\small] at (0,-1.95) {ningún $3$-coloreo mata las $168$ colineaciones; este $4$-coloreo sí};
  \end{tikzpicture}
  \caption{El plano de Fano ($7$ puntos, $7$ rectas), cuyas $168$ colineaciones son
  $\GL_3(\mathbb{F}_2)$ en los vectores no nulos de $\mathbb{F}_2^3$. Los colores
  únicos $1,2,3$ fijan los puntos $4,5,6$; tres rectas fijan entonces $2,0,1$ y la
  inyectividad fija $3$. Por tanto $D=4$.}
  \label{fig:fano}
\end{figure}

% ---------------------------------------------------------------
\section{Por qué importa}
\label{sec:why}

Me dirijo a tres lectores directamente, porque el trabajo habla a cada uno de forma
distinta.

\paragraph{A un teórico de grupos de permutaciones.} El caso extremal de la separación
suele encontrarse como una curiosidad en la frontera de un teorema de umbral. Leerlo
como una factorización exacta, y luego como una teselación, lo conecta con un cuerpo de
trabajo amplio y activo---los cánones aperiódicos de Vuza~\cite{Vuza1991}, las
condiciones de Coven--Meyerowitz~\cite{CovenMeyerowitz1999}, la pregunta de conjuntos
espectrales de Fuglede~\cite{Fuglede1974}---del que la literatura de sincronización no
se nutre ahora. El diccionario corre en ambos sentidos: construcciones de teselaciones
aperiódicas se vuelven construcciones de pares extremales no separables, y viceversa.
La forma órbita-local de la dicotomía es además, creo, el enunciado natural---no
necesita ni transitividad ni una cota de órbita global.

\paragraph{A quien construye matemáticas en Lean.} Estos cimientos no estaban antes
disponibles en un asistente de pruebas. La identidad de promedios, los teoremas de
uniformidad, la cota de base $D\le b+1$ y la API del número distinguidor son ahora
reutilizables; son exactamente las piezas sobre las que se construye la teoría de
matroides (IBIS) de la Parte~II. Mathlib gana un rincón pequeño pero portante de la
combinatoria de grupos de permutaciones que no tenía.

\paragraph{A quien trabaja en teoría de sincronización.} La jerarquía
separador/sincronizador se define precisamente a través de la condición de intersección
estudiada aquí. Hacer formal el motor de conteo y la dicotomía extremal, y atar el caso
extremal a una teoría de teselación bien desarrollada, ofrece al campo tanto un cimiento
comprobado como una fuente de ejemplos y contraejemplos que no ha estado explotando.
Los cánones de $\Z/n$ son un suministro inagotable de pares extremales.

% ---------------------------------------------------------------
\section{Dónde me detuve}
\label{sec:future}

Este es el primer conglomerado de una serie, honestamente mapeado, no una teoría
acabada. Hecho: separación de Neumann (finita e infinita); el motor de conteo y los
teoremas de uniformidad; la dicotomía órbita-local con la factorización y la
identificación con el canon rítmico; los números distinguidores, su API y la excepción
de Fano. Por delante, fijado pero aún no construido: el teorema IBIS de
Cameron--Fon-Der-Flaass~\cite{CFF1995} (las bases irredundantes de tamaño constante son
un matroide), el ladrillo central de la Parte~II y el hogar natural de la cota de base
probada aquí; los lemas puente explícitos que convierten este conglomerado de una lista
en un tejido; y los valores distinguidores restantes de Chan, junto con el análogo
proyectivo aparentemente no tratado $D(\mathrm{PGL}_n(\mathbb{F}_q)\curvearrowright
\mathrm{PG}^{n-1})$. El camino es clásico, pero no está construido.

% ---------------------------------------------------------------
\section{Trabajo relacionado}
\label{sec:related}

La matemática es clásica: el lema de separación es Birch--Burns--Oates
Macdonald--Neumann~\cite{BBMN1976} y la versión infinita de Neumann; la jerarquía de
conteo y la dicotomía separador/sincronizador son de Cameron y de
Araújo--Cameron--Steinberg~\cite{ACS2017}; la teoría extremal del caso $A=B$ es de
Barbieri et~al.~\cite{Barbieri2025}; los números distinguidores son de
Albertson--Collins y, para grupos lineales, de Chan~\cite{Chan2006} y
Klavžar--Wong--Zhu~\cite{KWZ2006}; el lado de la teselación es de Vuza y
Coven--Meyerowitz. En el lado formal, Mathlib aporta acciones de grupo,
órbita--estabilizador, Burnside, matroides y el lema de recubrimiento por cosets de
B.~H. Neumann, pero ninguno del lema de separación, la jerarquía de sincronización ni
los números distinguidores; no conozco tratamiento previo verificado por máquina de
estos en ningún asistente.

% ---------------------------------------------------------------
\section{Conclusión}

Una identidad de conteo, vista a través de distintas lentes, genera separación,
sincronización, factorizaciones exactas, cánones rítmicos y---por la pregunta dual del
coloreo---ruptura de simetría. La formalización no crea esta unidad; la hace explícita,
y comprobable, y convierte el cruce entre grupos de permutaciones y cánones musicales
de una analogía en un teorema. Nada de la matemática es nuevo; la contribución es que
ahora vive en un asistente de pruebas, y que el motor compartido es visible en el grafo
de dependencias en vez de afirmado. Cada pieza nueva surgió de mirar el caso que un
teorema clásico deja de lado---la igualdad en una desigualdad estricta, la frontera de
un umbral---que es a menudo donde se sienta un objeto sin nombre. Lo ofrezco como una
perla, y como la primera parte de una serie.

La Parte~II retoma el único ladrillo que este artículo solo nombra: el teorema de
Cameron y Fon-Der-Flaass de que un grupo de permutaciones cuyas bases irredundantes
tienen todas el mismo tamaño porta un matroide sobre sus puntos. Es el corazón
estructural al que ya apunta la cota de base $D\le b+1$, y con él las piezas reunidas
aquí---separación, conteo, uniformidad, distinción---se vuelven un capítulo conectado
en vez de una primera piedra. La identidad de conteo de la Sección~\ref{sec:intro}
seguirá por debajo; el objetivo de la serie es seguir haciendo visibles esos motores
compartidos, un ladrillo comprobado cada vez.

\paragraph{Sobre el método.} Este desarrollo se llevó a cabo con asistencia de IA
(Claude, Anthropic): definiciones, lemas y tácticas se propusieron y luego se aceptaron
o rechazaron contra el núcleo. El cuidado de alcance de este artículo---los núcleos
clásicos nombrados como tales, el único valor que no pude probar de forma factible en
Lean (el orden $168$) señalado y contrastado en GAP en vez de afirmado---refleja que el
núcleo, no el asistente, tiene la última palabra.

\paragraph{Artefactos.} Todas las fuentes Lean (\lean{sorry}-free; el informe de
axiomas está en la Tabla~\ref{tab:results}), los contrastes GAP de la
Tabla~\ref{tab:checks} y las figuras están en
\url{https://github.com/karlesmarin/godsil-gutman-lean}, archivado en
\textsc{doi}~\href{https://doi.org/10.5281/zenodo.21419285}{10.5281/zenodo.21419285}.

% ---------------------------------------------------------------
\appendix
\section{Los resultados formales}
\label{app:results}

La Tabla~\ref{tab:results} lista cada enunciado formalizado con su nombre Lean y su
estado. Todos los teoremas son \lean{sorry}-free y no vacíos, y \lean{\#print axioms}
reporta solo \lean{propext}, \lean{Classical.choice}, \lean{Quot.sound} (el
recubrimiento exhaustivo \lean{fanoWitnesses\_cover} necesita aún menos). Las tres
filas \stfuture{} son la frontera mapeada en la Sección~\ref{sec:future}.

\begin{table}[ht]
  \centering
  \small
  \arrayrulecolor{sepblue}
  \caption{Los resultados formalizados (Lean~4 / Mathlib).}
  \label{tab:results}
  \renewcommand{\arraystretch}{1.15}
  \begin{tabular}{@{}lll@{}}
    \toprule
    \rowcolor{sepblue}
    \textcolor{white}{\textbf{Nombre Lean}} & \textcolor{white}{\textbf{Enunciado}} & \textcolor{white}{\textbf{Estado}} \\
    \midrule
    \rowcolor{sepband}
    \texttt{separation\_lemma\_of\_pretransitive} & $|A||B|<|\Omega|\Rightarrow\exists g,\ gA\cap B=\emptyset$ & \stproven \\
    \texttt{neumann\_separation\_lemma} & órbitas de $A$ infinitas $\Rightarrow$ separable & \stproven \\
    \rowcolor{sepband}
    \texttt{sum\_card\_smul\_inter\_mul\_card} & $(\sum_g|gA\cap B|)\,|\Omega|=|A||B||G|$ \; (Cameron Thm~10) & \stproven \\
    \texttt{card\_smul\_inter\_eq\_of\_forall\_le} & promedio-igual-a-extremo \; (Cameron Thm~8) & \stproven \\
    \rowcolor{sepband}
    \texttt{IsSectionRegular.card\_part\_mul\dots} & sección-regular $\Rightarrow$ uniforme \; (Cameron Thm~9) & \stproven \\
    \texttt{separation\_dichotomy} & separa, \emph{o} $|gA\cap B|{=}1\ \forall g$ en una $|A||B|$-órbita & \stproven \\
    \rowcolor{sepband}
    \texttt{\dots\_eq\_one\_iff\_factorization} & extremal $\iff$ factorización exacta $G=B\cdot A^{-1}$ & \stproven \\
    \texttt{\dots\_vadd\_inter\_eq\_one\_iff\_tiling} & sobre $\Z/n$: extremal $\iff$ canon rítmico & \stproven \\
    \rowcolor{sepband}
    \texttt{distinguishingNumber} \small(+API) & mínimo $k$ con un $k$-coloreo que mata la simetría & \stdef \\
    \texttt{distinguishingNumber\_le\_card\_base\dots} & base de tamaño $b\Rightarrow D\le b{+}1$ \; (puente a IBIS) & \stproven \\
    \rowcolor{sepband}
    \texttt{fano\_distinguishingNumber} & $D(\GL_3(\mathbb{F}_2)\curvearrowright\mathbb{F}_2^3)=4$ & \stproven \\
    \midrule
    \emph{IBIS} (Cameron--Fon-Der-Flaass) & bases irredundantes tamaño cte.\ $\iff$ matroide & \stfuture \\
    \emph{lemas puente} & separación $\leftarrow$ motor; extremal $=$ Thm~8 $\ell{=}1$ & \stfuture \\
    \emph{clasificación completa de Chan} & los casos $D{=}3$; $D{=}2$ general; análogo $\mathrm{PGL}$ & \stfuture \\
    \bottomrule
  \end{tabular}
  \arrayrulecolor{black}
\end{table}

\bibliographystyle{plain}
\begin{thebibliography}{99}
\bibitem{BBMN1976} B.~J. Birch, R.~G. Burns, S. Oates Macdonald, P.~M. Neumann,
  \emph{On the orbit-sizes of permutation groups containing elements separating
  finite subsets}, Bull. Austral. Math. Soc. \textbf{14} (1976) 7--10.
\bibitem{BHNeumann1954} B.~H. Neumann, \emph{Groups covered by permutable
  subsets}, J. London Math. Soc. \textbf{29} (1954) 236--248.
\bibitem{ACS2017} J. Araújo, P.~J. Cameron, B. Steinberg, \emph{Between primitive
  and 2-transitive: synchronization and its friends}, EMS Surv. Math. Sci.
  \textbf{4} (2017) 101--184.
\bibitem{CameronSync} P.~J. Cameron, \emph{Synchronization}, lecture notes
  (LTCC), 2010; \url{https://cameroncounts.wordpress.com/lecture-notes/}.
\bibitem{Barbieri2025} M. Barbieri, P. Lekše, P. Potočnik, K. Rekvényi,
  \emph{Separating subsets from their images} (2025), arXiv:2508.20731.
\bibitem{CFF1995} P.~J. Cameron, D.~G. Fon-Der-Flaass, \emph{Bases for
  permutation groups and matroids}, European J. Combin. \textbf{16} (1995)
  537--544.
\bibitem{AlbertsonCollins1996} M.~O. Albertson, K.~L. Collins, \emph{Symmetry
  breaking in graphs}, Electron. J. Combin. \textbf{3} (1996) \#R18.
\bibitem{Tymoczko2004} J. Tymoczko, \emph{Distinguishing numbers for graphs and
  groups}, Electron. J. Combin. \textbf{11} (2004) \#R63.
\bibitem{Chan2006} M. Chan, \emph{The maximum distinguishing number of a group},
  Electron. J. Combin. \textbf{13} (2006) \#R70.
\bibitem{KWZ2006} S. Klavžar, T.-L. Wong, X. Zhu, \emph{Distinguishing labelings
  of group action on vector spaces and graphs}, J. Algebra \textbf{303} (2006)
  626--641.
\bibitem{CovenMeyerowitz1999} E.~M. Coven, A. Meyerowitz, \emph{Tiling the
  integers with translates of one finite set}, J. Algebra \textbf{212} (1999)
  161--174.
\bibitem{Vuza1991} D.~T. Vuza, \emph{Supplementary sets and regular complementary
  unending canons (Parts One--Four)}, Perspectives of New Music \textbf{29}(2)
  (1991), \textbf{30}(1),(2) (1992), \textbf{31}(1) (1993).
\bibitem{Fuglede1974} B. Fuglede, \emph{Commutative sets of self-adjoint partial
  differential operators and the group of orthogonal exponentials}, J. Funct.
  Anal. \textbf{16} (1974) 101--121.
\bibitem{Amiot2016} E. Amiot, \emph{Music Through Fourier Space: Discrete Fourier
  Transform in Music Theory}, Springer, 2016.
\bibitem{Mathlib2020} The mathlib Community, \emph{The Lean mathematical library},
  CPP 2020.
\end{thebibliography}

\end{document}
